\section{Experimental basis and input reconstruction}
\label{sec:experiments}
\subsection{Available configurations}
The experimental columns of Liu et al.\ \cite{Liu2025chemical} were
1.25 m long, with an inner diameter of 4.8 mm for 316L stainless steel
and 5.0 mm for fused silica. A 15 mm sample boat was inserted into a
hot starting furnace, followed by a furnace providing a decreasing
temperature profile. Helium carried the evaporated material downstream.
The source was withdrawn after three hours; the furnaces and gas flow
were then turned off and the column cooled before deposit analysis.
Temperature measurements and radiotracer scans used 1 cm increments.

\Cref{tab:experiments} lists the three configurations for which both
temperature curves and iodine scans could be reconstructed. The
stand-alone BiI$_3$ experiment at 800\,$^\circ$C was considered in case
preparation but was not run, because its required axial temperature
profile was not available in the supplied figures. The missing case is
not assigned a surrogate profile or counted as an unsuccessful simulation.
\begin{table}[tb]
\centering
\caption{Published configurations used for the baseline matrix. Masses
refer to the initial sample, not to the amount released in a 60 s run.
All experimental evaporation periods were 10,800 s.}
\label{tab:experiments}
\begin{tabular}{llrrr}
\toprule
Identifier & Column & Diameter (mm) & Sample (mg) & He (mL/min) \\
\midrule
PbI2\_SS & Steel & 4.8 & 14.6 & 100 \\
LBE-I\_SS\_II & Steel & 4.8 & 346.0 & 45 \\
LBE-I\_SiO2\_I & Silica & 5.0 & 368.3 & 45 \\
\bottomrule
\end{tabular}
\end{table}

The two LBE samples had reported iodine mole fraction
$9.04\times10^{-4}$ and a 900\,$^\circ$C starting furnace. Their
reported deposited iodine amounts were 0.191 and 0.203 mg, respectively.
These deposited amounts constrain the prescribed mean iodine source in
the present cases; they are not direct measurements of a time-resolved
evaporation rate or an independent inlet-speciation measurement.

\subsection{Digitisation and coordinate assumptions}
Axial temperature curves and iodine histograms were digitised from the
published figures. Temperatures are linearly interpolated at cell centres
and wall faces; values beyond the plotted interval are held at the nearest
endpoint. The computational coordinate is identified with the plotted
coordinate. Its origin relative to the physical boat and furnace has not
been independently confirmed. \Cref{fig:temperature} displays the
reconstructed curves. These assumptions are retained in the original
baseline matrix and varied explicitly in the silica-coordinate study.

The recorded digitisation resolutions are approximately 1 mm in position,
4 K in temperature and 0.4 percentage points in histogram height. They
describe raster extraction, not the full experimental uncertainty. They
do not include coordinate registration, detector response, occluded curve
segments, temperature reproducibility or calibration. Likewise, the
uncertainties printed alongside experimental deposition temperatures
must not be interpreted as uncertainties on every digitised profile bin.
The reconstructed histogram sums are 99.17\%, 96.70\% and 102.24\%
for pure PbI$_2$, steel LBE and silica LBE. The primary comparison keeps
those values without silently renormalising them.

In the LBE scans, the paper assigns 6.6\% and 7.0\% of iodine to KI,
which is outside the Pb--Bi--I species set. The model profile is scaled
to the remaining 93.4\% or 93.0\%, while the measured histogram retains
all its published deposits. A second, explicitly conditional comparison
normalises both profiles over their common scanned range
(\cref{sec:comparison}). Neither comparison supplies the omitted
potassium chemistry.

\begin{figure}[tb]
\centering
\includegraphics[width=.88\textwidth]{temperature_profiles.pdf}
\caption{Temperature curves digitised from Liu et al.\
\cite{Liu2025chemical}, with the assumed computational source interval.
Unplotted parts of the CFD domain use constant endpoint extensions.
The curves are experimental figure reconstructions, not temperature
fields predicted by a conjugate heat-transfer calculation.}
\label{fig:temperature}
\end{figure}

\subsection{Simulation duration and measurement protocol}
All source-driven cases run for 60 s while retaining a source rate based
on the 10,800 s experimental period. Thus they release $1/180$ of the
assumed three-hour inventory; the experimental inventory is not compressed
into a one-minute calculation. Continuous-source stationarity can inform
deposition rates and conditional shapes, but cumulative wall inventories
retain the initial filling transient. Four additional cases stop the
source and carrier flow at 60 s and continue to 80 s under the original
temperature profile. This probes residual-gas relaxation, not the
experimental cooling history.
